% Источники: 02_Общие_соглашения/Методика_идентификаторов.md
% Генератор: docs/build_pdf.py
% Шаблон печатного представления. Содержимое подставляет docs/build_pdf.py.
\documentclass[12pt,a4paper]{article}
\usepackage{iftex}
\ifPDFTeX
  \usepackage[utf8]{inputenc}
  \usepackage[T2A]{fontenc}
  \usepackage{newtxtext,newtxmath}
\else
  \usepackage{fontspec}
  \IfFontExistsTF{Times New Roman}
    {\setmainfont{Times New Roman}}
    {\setmainfont{Libertinus Serif}}
  \IfFontExistsTF{Arial}{\setsansfont{Arial}}{\setsansfont{Latin Modern Sans}}
  \IfFontExistsTF{Menlo}{\setmonofont{Menlo}[Scale=MatchLowercase]}
    {\setmonofont{Latin Modern Mono}[Scale=MatchLowercase]}
\fi
\usepackage[russian]{babel}
\usepackage{geometry}
\geometry{left=25mm,right=15mm,top=20mm,bottom=20mm}
\usepackage{setspace}
\onehalfspacing
\usepackage{indentfirst}
\setlength{\parindent}{1.25cm}
\setlength{\parskip}{2pt}
\setlength{\emergencystretch}{3em}
\usepackage{graphicx,booktabs,longtable,array,tabularx,calc,etoolbox}
\usepackage{enumitem}
\setlist[itemize]{leftmargin=1.25cm,itemsep=2pt,topsep=4pt}
\setlist[enumerate]{leftmargin=1.25cm,itemsep=2pt,topsep=4pt}
\newcommand{\tightlist}{\setlength{\itemsep}{0pt}\setlength{\parskip}{0pt}}
\AtBeginEnvironment{longtable}{\small\singlespacing}
\setlength{\LTpre}{8pt}
\setlength{\LTpost}{8pt}
\setlength{\tabcolsep}{5pt}
\setlength{\LTleft}{0pt}
\setlength{\LTright}{0pt}
\usepackage{xurl}
\usepackage[unicode,colorlinks=true,linkcolor=black,urlcolor=blue]{hyperref}
\usepackage{bookmark}
\hypersetup{pdftitle={Методика идентификаторов},pdfauthor={Анохин Г.А.},pdfsubject={Криптотренажёр mephi-12}}
\urlstyle{same}
\pagestyle{plain}
\setcounter{secnumdepth}{3}
\setcounter{tocdepth}{2}

\begin{document}
\hypersetup{pageanchor=false}
\begin{titlepage}
  \thispagestyle{empty}
  \centering
  {\large НИЯУ МИФИ\par}
  \vspace*{3.5cm}
  {\LARGE\bfseries Методика идентификаторов\par}
  \vspace{0.5cm}
  {\large Методика ID\par}
  \vspace{0.6cm}
  {\large Платформа для изучения и практики курса по криптографии\par}
  {\large Криптотренажёр mephi-12\par}
  \vspace{0.7cm}
  {\large Версия 0.3\par}
  {\large Редакция от 07.11.2026\par}
  \vspace{1cm}
  \begin{tabularx}{\textwidth}{@{}p{3.4cm}Xp{4cm}@{}}
    \toprule
    Роль & ФИО & Подпись / Дата \\
    \midrule
    Разработал & Анохин Г.А. & \\
    \addlinespace[6pt]
    Согласовано & & \\
    \addlinespace[6pt]
    Утверждено & & \\
    \bottomrule
  \end{tabularx}
  \vfill
  {\large Москва, 2026\par}
\end{titlepage}
\hypersetup{pageanchor=true}

\section*{Лист регистрации изменений}
\addcontentsline{toc}{section}{Лист регистрации изменений}
\begingroup\small\singlespacing
\noindent
\begin{tabularx}{\textwidth}{@{}p{1.2cm}p{2cm}p{2.5cm}Xp{1.8cm}@{}}
\toprule
Версия & Дата & Автор & Содержание изменения & Утвердил \\
\midrule
0.3 & 07.11.2026 & Анохин Г.А. & Уточнены правила и контракты; расширено описание сценариев; подготовлена печатная редакция. & \\
\bottomrule
\end{tabularx}
\endgroup
\clearpage
\begingroup\small\singlespacing
\tableofcontents
\endgroup
\clearpage

\hypertarget{методика-идентификаторов}{}

Версия: 0.3. Дата: 07.11.2026. Ответственный: Анохин Г.А. Методика действует для требований этой редакции.

\section{Область действия}\label{ux43eux431ux43bux430ux441ux442ux44c-ux434ux435ux439ux441ux442ux432ux438ux44f}

Методика задаёт обозначения требований и связей в \href{https://github.com/Kokorev-VD/mephi-samolety/blob/main/docs/README.md}{текущей редакции}. Схема «тип-домен-номер» сохраняет идентификаторы ранее подготовленных Markdown-требований сессий. Это соглашение проекта; оно не объявляется обязательным форматом какого-либо стандарта.

В исходной \href{https://github.com/Kokorev-VD/mephi-samolety/blob/main/notes/methodology/id-methodology.md}{методике комплекта} используются плоские номера \texttt{SR-001}, \texttt{SWR-001}. Они не переименовываются этой редакцией. При переносе редактор указывает соответствие между прежним и новым обозначением; оно должно подтверждаться смыслом требования. Совпадение номера не доказывает совпадение смысла.

\section{Форматы}\label{ux444ux43eux440ux43cux430ux442ux44b}

{\def\LTcaptype{none} % do not increment counter
\begin{longtable}[]{@{}
  >{\raggedright\arraybackslash}p{(\linewidth - 2\tabcolsep) * \real{0.3200}}
  >{\raggedright\arraybackslash}p{(\linewidth - 2\tabcolsep) * \real{0.6800}}@{}}
\toprule\noalign{}
\begin{minipage}[b]{\linewidth}\raggedright
Объект
\end{minipage} & \begin{minipage}[b]{\linewidth}\raggedright
Обозначение
\end{minipage} \\
\midrule\noalign{}
\endhead
\bottomrule\noalign{}
\endlastfoot
Документ & \texttt{DOC-NN}; для подсистем допускается суффикс \texttt{.N} \\
Системное требование & \texttt{SR-ДОМЕН-NNN} \\
Программное требование & \texttt{SWR-ДОМЕН-NNN} \\
Описательный проверочный кейс & \texttt{TC-ДОМЕН-NNN} \\
Зарегистрированное замечание & \texttt{DEF-NNN} \\
\end{longtable}
}

Номер начинается с \texttt{001}; для SR, SWR и TC он выдаётся отдельно внутри каждого домена. В обозначении используются латинские буквы, цифры и обычный дефис. Маркеры требований действующего ТЗ, например \texttt{[TS.\allowbreak{}REQ.\allowbreak{}FUNC.\allowbreak{}SESSION.\allowbreak{}CREATE]}, сохраняются в исходном виде.

{\def\LTcaptype{none} % do not increment counter
\begin{longtable}[]{@{}
  >{\raggedright\arraybackslash}p{(\linewidth - 4\tabcolsep) * \real{0.3219}}
  >{\raggedright\arraybackslash}p{(\linewidth - 4\tabcolsep) * \real{0.3563}}
  >{\raggedright\arraybackslash}p{(\linewidth - 4\tabcolsep) * \real{0.3219}}@{}}
\toprule\noalign{}
\begin{minipage}[b]{\linewidth}\raggedright
Домен
\end{minipage} & \begin{minipage}[b]{\linewidth}\raggedright
Область
\end{minipage} & \begin{minipage}[b]{\linewidth}\raggedright
Владелец по плану
\end{minipage} \\
\midrule\noalign{}
\endhead
\bottomrule\noalign{}
\endlastfoot
\texttt{AUTH} & Аутентификация и пользователи & Кокорев \\
\texttt{SES} & Сессии задач и планировщик & Анохин \\
\texttt{KNAP} & Ранцевые криптосистемы & Георгиев \\
\texttt{BDR} & Учебная бэкдор-атака & Минасов \\
\texttt{ELG} & Эль-Гамаль & Портнов \\
\texttt{CFG} & Конфигурации, модерация и инфраструктура & Самбьюи \\
\texttt{SYS} & Общие требования и интеграция & Редактор совместно с владельцами \\
\end{longtable}
}

\section{Правила ведения}\label{ux43fux440ux430ux432ux438ux43bux430-ux432ux435ux434ux435ux43dux438ux44f}

Автор выдаёт свободный ID в своём домене. Один ID обозначает одну проверяемую обязанность или неделимую реакцию на событие. При уточнении того же требования ID сохраняется. Для замены другим требованием указываются старый и новый ID; отменённый номер повторно не выдаётся. Перенос между файлами не меняет ID.

Каждое требование содержит формулировку, источник, приоритет и способ проверки. Общие значения этих полей можно задать перед группой требований. Приоритеты: обязательный, желательный, необязательный. Приоритет требования не заменяет состояние документа или факт его согласования.

В источнике указывается требование ТЗ, более высокого уровня или предметный сценарий. Если исходные документы расходятся, связь сопровождается пояснением; её нельзя выдавать за полное соответствие. Общее требование определяется один раз. В родительских документах и таблицах используются ссылки на его определение.

\section{Трассировка}\label{ux442ux440ux430ux441ux441ux438ux440ux43eux432ux43aux430}

Для подготовленных программных требований указываются реальные системные источники. Связи с TC добавляются после подготовки самих кейсов. До этого в таблице указывается описательный способ проверки без вымышленного номера TC.

Полную матрицу SR, SWR и TC ведёт Портнов. Автор домена передаёт ему свои определения и связи. Назначение номера кейсу не означает, что проверка выполнена. DEF присваивается только зарегистрированному замечанию фактической инспекции.

История: 05.11.2026 --- подготовлена рабочая редакция для второго спринта; сохранены \texttt{SR-SES-001}--\texttt{SR-SES-010}, добавлено правило учёта времени поступления ответа.


\end{document}
